% Generated draft: requires mathematical and source review.
\documentclass[11pt]{article}
\usepackage[margin=25mm]{geometry}
\usepackage{fontspec}
\setmainfont{Latin Modern Roman}
\usepackage{amsmath,amssymb,mathtools,unicode-math}
\setmathfont{Latin Modern Math}
\usepackage{enumitem,needspace,xurl,hyperref}
\setlength{\parindent}{0pt}
\setlength{\parskip}{5pt}
\setlength{\emergencystretch}{3em}
\newcommand{\N}{\mathbb{N}}
\newcommand{\Z}{\mathbb{Z}}
\newcommand{\Q}{\mathbb{Q}}
\newcommand{\R}{\mathbb{R}}
\newcommand{\C}{\mathbb{C}}
\newcommand{\F}{\mathbb{F}}
\newcommand{\A}{\mathbb{A}}
\newcommand{\PP}{\mathbb{P}}
\newcommand{\E}{\mathbb{E}}
\newcommand{\eps}{\varepsilon}
\newcommand{\dd}{\,\mathrm{d}}
\begin{document}
\section*{UnsolvedMath ID: 20001272}
\textbf{Source problem code:} catalogue entry
\subsection{Definitions and mathematical statement}
Fix \(n\ge3\) and \(2\le k\le n-1\), as specified by the source introduction, and let \(G(n,k)\) be the Grassmannian of \(k\)-dimensional linear subspaces of \(\R^n\). The supplied verification repairs ``for every \(G(n,k)\)'' to ``for every \(S\in G(n,k)\).'' The question asks whether corresponding homothetic central sections of star bodies force the bodies themselves to be homothetic.

For the aid's explicit direct-homothety setting, take origin-star bodies
\[
 K=\{tu:u\in\mathbb S^{n-1},\ 0\le t\le\rho_K(u)\},
\]
and similarly \(L\), with positive continuous radial functions. Assume
\[
 \forall S\in G(n,k)\ \ \exists a(S)\in S,\ \lambda(S)>0:
 K\cap S=a(S)+\lambda(S)(L\cap S).
\]
The direct version of the source question asks whether \(K=a+\lambda L\) for some \(a\in\R^n\), \(\lambda>0\). The original excerpt does not define star-body center conventions or whether negative homothety ratios are allowed, so the aid's positive-ratio model is explicit. The source reports an affirmative answer when the bodies are convex and \(o\in\operatorname{int}K\cap\operatorname{int}L\), citing [G, Theorem 7.1.1].

For a section \(A\subset S\), write \(V_k(A)\) for its Euclidean \(k\)-volume and \(g(A)=V_k(A)^{-1}\int_Ax\,\dd x\) for its centroid. The aid reports uniquely determined, continuous parameters
\[
 \lambda(S)=\left(\frac{V_k(K\cap S)}{V_k(L\cap S)}\right)^{1/k},
 \qquad
 a(S)=g(K\cap S)-\lambda(S)g(L\cap S).
\]
Continuity is on the Grassmannian, viewing \(a(S)\in S\) as a vector in \(\R^n\).

Its rigidity result is: if \(a(S)=0\) for every \(S\), then there is one constant \(\lambda>0\) with \(\rho_K(u)/\rho_L(u)=\lambda\) for every unit direction \(u\), hence \(K=\lambda L\). The hypothesis holds when both corresponding section centroids are zero, and in particular when both star bodies are origin-symmetric. No convexity is needed for these special cases. The restriction \(k\ge2\) lets a common section contain any two directions and is retained.

The aid explicitly warns that general translated homotheties cannot automatically be restricted to intersections of section planes: translation by \(a(S)\) need not preserve such an intersection. It does not prove that every canonical translation field vanishes or comes from one global homothety. Thus continuity and uniqueness of the section parameters alone do not settle the general nonconvex question.
\subsection{Short English statement}
The source asks whether homothetic central sections determine star bodies up to homothety. The supplied aid gives canonical continuous section parameters and proves global dilation when their translations vanish, covering centered-section and origin-symmetric cases without convexity.
\subsection{Sources}
\begin{itemize}
\item Local UnsolvedMath registry, permanent ID \(20001272\), source code \texttt{catalogue entry}; \texttt{aim-convex-geometry-notes.json}, \texttt{notes[3]}, tag \texttt{problem}.
\item Richard Gardner's problem list for the 2013 AIM workshop ``Sections of Convex Bodies,'' Problem 3: \url{https://aimath.org/pastworkshops/sectionsconvexproblems.pdf}. Statement status \texttt{exact}; the supplied verification confirms the introductory \(k\)-range, Grassmannian notation, and missing \(S\in\) repair. No PDF was accessed for this draft.
\item Source labels [G, Problem 7.1 and Note 7.1] and [G, Theorem 7.1.1], the latter for convex bodies with the origin in both interiors. The bibliographic entry [G] is not supplied and is not invented.
\item Research attempt 1: \texttt{partial\_result}, problem status at run \texttt{open}, rigor label \texttt{proved}, novelty confidence low; registry status \texttt{partially\_solved}, classification \texttt{PARTIAL-PROGRESS}. Created and updated 14 August 2026. No original remarks or literature field are supplied. This machine-generated aid is not peer reviewed; independently verify before citation.
\item Ulam AI / UnsolvedMath Contributors; CC BY 4.0.
\end{itemize}
\end{document}
